\section{Experiments}
\label{sec:experiments}

\subsection{Setup}
\label{sec:setup}

\textbf{Backbone.}
All experiments use ViT-B/16 pretrained on ImageNet-21K,
with the backbone frozen throughout training.
LoRA adapters are inserted into the query and value projections
of all 12 transformer blocks (24 adapter pairs total).

\textbf{Evaluation tracks.}
We report results on two complementary benchmark tracks:

\emph{Track~1 — Fine-grained recognition (6 datasets, $r=16$).}
These datasets span diverse visual categories and task structures:
CUB200 (10t $\times$ 20c),
Cars196 (10t $\times$ 20c),
Aircraft (10t $\times$ 10c),
Flowers (10t $\times$ 10c),
OxfordPet (9t $\times$ 4c),
ImageNet-R (20t $\times$ 10c).
This track directly tests the orthogonality principle across varied domains.

\emph{Track~2 — Standard PECL benchmarks (4 datasets, $r=10$--$30$).}
Following the protocol of~\citet{zheng2025ewclora}:
CIFAR-100 (10t $\times$ 10c, $r=10$),
ImageNet-A (10t $\times$ 20c, $r=10$),
ImageNet-R (10t $\times$ 20c, $r=10$),
DomainNet (5t $\times$ 69c, $r=30$).
This track provides direct comparability against the EWC-LoRA baseline.

\textbf{Metrics.}
(1)~\textbf{CNN FAA}: Final Average Accuracy (mean over all tasks at end of training);
(2)~\textbf{CNN AAA}: Average Anytime Accuracy (lower-triangle mean, measures retention throughout);
(3)~\textbf{Forget}: mean per-task forgetting;
(4)~\textbf{BWT}: backward transfer.
All metrics use the cosine-normalized classification head (CNN).

\textbf{Optimizer.}
We use SGD as the base optimizer for GAM, with cosine learning rate schedule.
Track~1: $\mathrm{lr}=0.01$, $\rho=0.05$, gradient-norm perturbation radius $0.1$, 50 epochs for ImageNet-R and 20 epochs otherwise.
Track~2: $\mathrm{lr}=0.01$--$0.05$ (per-dataset tuned; see Appendix~\ref{sec:impl}).
Fixed across both tracks: $\lambdaflat=1.0$, $\lambdaewc=2000$, $\gamma=0.9$,
Fisher over 100 mini-batches at task boundary.

\textbf{Reproducibility note.}
All Track~1 results reported in the main tables are single-seed runs (seed~0) due to compute constraints.
The reported improvements of $\lesssim 1$pp on CUB200 should be interpreted as indicative
rather than statistically certified; we flag this as a limitation.
Multi-seed evaluation is planned for the final submission.

\textbf{Method categories.}
We distinguish two categories of baselines:
\begin{itemize}
  \item \textbf{Single-adapter methods} maintain a fixed set of LoRA parameters shared across all tasks.
    This includes SeqLoRA (no protection), EWC-LoRA (Fisher regularization only),
    and FS-LoRA (flatness + Fisher, ours).
  \item \textbf{Structural methods} grow or isolate parameter capacity per task
    (InfLoRA, SDLoRA, IncLoRA, O-LoRA).
    These are shown for context but are not directly comparable to single-adapter methods.
\end{itemize}
All baselines are run with GAM optimizer to ensure a fair comparison of the \emph{regularization design},
independent of the base optimizer.

\subsection{Main Results: Fine-Grained Recognition}
\label{sec:main_results}

Table~\ref{tab:main} reports CNN FAA on all six Track~1 datasets.
We compare FS-LoRA against the single-adapter baseline SeqLoRA+GAM
and against four structural LoRA-CL methods run with the same GAM optimizer.

\begin{table}[t]
\caption{CNN FAA on Track~1 fine-grained recognition benchmarks.
All methods use GAM as the base optimizer and ViT-B/16 frozen backbone, no replay.
Single-adapter group uses a fixed shared LoRA ($r=16$) throughout.
Structural methods grow or isolate parameters per task (shown for context; marked~$\dagger$).
Best single-adapter result per dataset in \textbf{bold};
overall best (including structural) underlined.}
\label{tab:main}
\centering
\small
\begin{tabular}{llcccccc}
\toprule
Category & Method & CUB200 & Cars196 & Aircraft & Flowers & OxfordPet & ImageNet-R \\
\midrule
\multirow{3}{*}{\shortstack[l]{Single\\adapter}}
  & SeqLoRA+GAM      & 73.249 & 52.197 & 47.081 & 81.282 & 80.921 & 72.986 \\
  & Layer 2 (decoupled) & 74.220 & 51.975 & 46.660 & ---    & ---    & 73.318 \\
  & \textbf{FS-LoRA (ours)} & \textbf{74.224} & \textbf{51.741} & \textbf{46.660} & \textbf{86.747} & 78.824 & \textbf{73.143} \\
\midrule
\multirow{4}{*}{\shortstack[l]{Structural\\($\dagger$)}}
  & InfLoRA+GAM$^\dagger$  & \underline{75.169} & 43.413 & 45.458 & \underline{87.652} & \underline{81.433} & --- \\
  & SDLoRA+GAM$^\dagger$   & 70.477 & 38.218 & 40.899 & 84.615 & 79.938 & --- \\
  & IncLoRA+GAM$^\dagger$  & 68.769 & \underline{53.353} & 46.060 & 81.134 & 79.873 & 72.943 \\
  & O-LoRA+GAM$^\dagger$   & 68.518 & 50.640 & 39.278 & 80.952 & 79.746 & 72.636 \\
\bottomrule
\end{tabular}
\end{table}

\noindent\small{ImageNet-R uses inc=10, 20 tasks, $r=16$.
IncLoRA 72.943, O-LoRA 72.636, both below FS-LoRA \textbf{73.143} (bold).
``---'' indicates no result available under the matching setting.}

\textbf{Single-adapter comparison.}
FS-LoRA shows mixed results relative to SeqLoRA+GAM:
CUB200 ($+1.0$pp), Cars196 ($-0.5$pp), Aircraft ($-0.4$pp), Flowers ($+5.5$pp),
OxfordPet ($-2.1$pp), ImageNet-R ($+0.2$pp).
Positive gains on CUB200, Flowers, and ImageNet-R; negative or negligible on Cars196,
Aircraft, and OxfordPet.
The largest gain is Flowers ($+5.5$pp), where fisher/clean $=0.010$, indicating
AS$^1$ flatness is the primary driver (Fisher penalty is lightly active).
Cars196 shows a small FAA loss but positive AAA gain ($+0.85$pp), suggesting improved
retention during training but a marginally lower final accuracy.
All Track~1 results are \emph{single-seed} runs; gains of $\leq 1$pp should be
interpreted as indicative, not statistically confirmed.
OxfordPet is an outlier with flat/clean $=204$; see Section~\ref{sec:discussion}.

\textbf{Comparison with structural methods.}
InfLoRA+GAM surpasses FS-LoRA on CUB200 ($75.2$ vs $74.2$) and Flowers ($87.7$ vs $86.7$),
using structural subspace allocation.
FS-LoRA matches or exceeds IncLoRA and O-LoRA on five of the six datasets available,
despite using a \emph{fixed} shared adapter.
This demonstrates that principled regularization design can close much of the gap
to structural approaches.

\subsection{Retention and Forgetting}
\label{sec:retention}

Table~\ref{tab:forget} reports CNN AAA and Forget alongside FAA for the same methods,
restricted to datasets where complete results are available for FS-LoRA.

\begin{table}[h]
\caption{CNN FAA / AAA / Forget on CUB200, Aircraft, and Flowers (complete runs).
FS-LoRA improves AAA over SeqLoRA+GAM on all three datasets,
confirming better retention throughout training.}
\label{tab:forget}
\centering
\small
\begin{tabular}{lcccccccccc}
\toprule
& \multicolumn{3}{c}{CUB200} & \multicolumn{3}{c}{Aircraft} & \multicolumn{3}{c}{Flowers} \\
\cmidrule(lr){2-4}\cmidrule(lr){5-7}\cmidrule(lr){8-10}
Method & FAA & AAA & Fgt & FAA & AAA & Fgt & FAA & AAA & Fgt \\
\midrule
SeqLoRA+GAM & 73.25 & 73.54 & 11.55 & 47.08 & 43.87 & 16.80 & 81.28 & 85.80 & 8.02 \\
InfLoRA+GAM$^\dagger$ & 75.17 & 76.19 & 13.28 & 45.46 & 42.15 & 22.07 & 87.65 & 90.15 & 3.18 \\
FS-LoRA (ours) & \textbf{74.22} & \textbf{74.33} & \textbf{11.71} & \textbf{46.66} & \textbf{43.86} & \textbf{15.57} & \textbf{86.75} & \textbf{89.74} & \textbf{4.58} \\
\bottomrule
\end{tabular}
\end{table}

\textbf{Key observation on forgetting.}
InfLoRA+GAM achieves lower forgetting on Flowers ($3.2$ vs $4.6$) using structural isolation,
but its forgetting on Aircraft ($22.1$) is \emph{higher} than SeqLoRA+GAM ($16.8$) and FS-LoRA ($15.6$).
This suggests that structural allocation alone does not guarantee stable retention,
whereas FS-LoRA's Fisher penalty provides consistent anchoring across domains.

\subsection{Ablation: Three-Layer Design on CUB200}
\label{sec:ablation}

Table~\ref{tab:ablation} traces the three-step design progression on CUB200,
isolating the contribution of decoupling and trace normalization.

\begin{table}[h]
\caption{Design ablation on CUB200 (10 tasks $\times$ 20 classes, seed 0).
Layer 1: GAM applied to $\Ltask+\Lfisher$ (coupled).
Layer 2: decoupled GAM on $\Ltask$ only, raw Fisher.
Layer 3 (FS-LoRA): decoupled + trace-normalized Fisher.
Fisher/clean = $\|\gfisher\|/\|\gclean\|$ (mean over tasks).}
\label{tab:ablation}
\centering
\begin{tabular}{llccccc}
\toprule
Layer & $\lambdaewc$ & Normalized & FAA & AAA & Forget & Fisher/clean \\
\midrule
1 (coupled) & 20   & No  & 73.227 & 72.952 & 10.796 & --- \\
2 (decoupled) & 20   & No  & 73.845 & 74.245 & 12.079 & 0.0029 \\
2 (decoupled) & 100  & No  & 74.182 & 74.309 & 11.806 & 0.0137 \\
2 (decoupled) & 1000 & No  & 74.220 & 74.284 & 11.563 & 0.0588 \\
\midrule
\textbf{3 / FS-LoRA} & \textbf{2000} & \textbf{Yes} & \textbf{74.224} & \textbf{74.332} & \textbf{11.709} & \textbf{0.0248} \\
3 / FS-LoRA & 5000 & Yes & 74.181 & 74.255 & 11.667 & --- \\
\bottomrule
\end{tabular}
\end{table}

\textbf{Note on missing EWC-only baseline.}
The ablation does not include a ``Fisher only, no flatness'' (SeqLoRA+EWC, no GAM) row,
because this baseline was not run on the fine-grained datasets.
SeqLoRA+SGD on CUB200 gives FAA$=57.8$ vs SeqLoRA+GAM $73.2$, confirming that
flatness alone provides a large gain;
the incremental gain from adding Fisher ($74.2$) is modest but consistent.
A full $2\times2$ ablation (flatness$\times$Fisher) is listed as future work.

\textbf{Observations:}
\begin{enumerate}
  \item \textbf{Decoupling matters.}
    Moving from Layer~1 to Layer~2 at $\lambdaewc=20$ improves FAA by $0.6$pp.
    The GAM closure in Layer~1 is distorted by the Fisher penalty,
    which inflates the update norm and reduces gradient informativeness.
  \item \textbf{Raw Fisher is nearly inactive.}
    At $\lambdaewc=20$ without normalization, Fisher/clean $=0.003$---the Fisher penalty
    contributes less than $0.3\%$ of the clean gradient magnitude and is effectively inactive.
    Performance gains at Layer~2 come almost entirely from AS$^1$ flatness.
  \item \textbf{Trace normalization restores Fisher activity.}
    With trace normalization, $\lambdaewc=2000$ achieves Fisher/clean $=0.025$---comparable
    to Layer~2 at $\lambdaewc=1000$---while using a dimensionless, dataset-agnostic hyperparameter.
    The Fisher component now genuinely contributes to the update.
  \item \textbf{Stability across $\lambdaewc$.}
    FS-LoRA FAA varies by less than $0.05$pp across $\lambdaewc\in\{2000,5000\}$,
    confirming that trace normalization provides consistent behavior.
\end{enumerate}

\subsection{Standard PECL Benchmarks}
\label{sec:paper4}

Table~\ref{tab:paper4} compares FS-LoRA against the EWC-LoRA reference results
from~\citet{zheng2025ewclora} on the four standard PECL benchmarks (Track~2).
ImageNet-A results are complete for two LR settings; other datasets are partial runs.

\begin{table}[h]
\caption{CNN FAA on standard PECL benchmarks (Track~2),
compared against EWC-LoRA~\citep{zheng2025ewclora}.
``$\dag$'' denotes a complete final run; ``(partial)'' indicates the run is in progress.
FS-LoRA uses SGD base optimizer with per-dataset tuned learning rate.}
\label{tab:paper4}
\centering
\begin{tabular}{lcccc}
\toprule
Method & CIFAR-100 & ImageNet-A & ImageNet-R & DomainNet \\
\midrule
EWC-LoRA~\citep{zheng2025ewclora} & 87.91 & 59.89 & 72.86 & 73.46 \\
\midrule
FS-LoRA (SGD, $\mathrm{lr}=0.01$) & --- & 45.999$^\dag$ & --- & --- \\
FS-LoRA (SGD, $\mathrm{init\_lr}=0.04$, $\mathrm{lr}=0.02$) & --- & \textbf{58.92}$^\dag$ & --- & --- \\
FS-LoRA (SGD, lr-tuned) & --- & --- & --- & --- \\
\midrule
\multicolumn{5}{l}{\emph{Partial runs (SGD, $\mathrm{lr}=0.01$, for reference only)}} \\
FS-LoRA (CIFAR-100, t7/10)  & (88.37) & --- & --- & --- \\
FS-LoRA (ImageNet-R, t9/10) & --- & --- & (77.13) & --- \\
FS-LoRA (DomainNet, t2/5)   & --- & --- & --- & (73.69) \\
\bottomrule
\end{tabular}
\end{table}

\textbf{ImageNet-A: LR is the critical factor.}
With $\mathrm{lr}=0.01$, FS-LoRA achieves FAA~$=45.999$ ($-13.9$pp vs.\ EWC-LoRA).
With a warm-start cosine schedule ($\mathrm{init\_lr}=0.04$, $\mathrm{lr}=0.02$),
FS-LoRA reaches FAA~$=\mathbf{58.92}$ ($-0.97$pp vs.\ EWC-LoRA), Forget~$=5.6$.
The NME-head FAA is $64.45$, surpassing the EWC-LoRA CNN reference.
This confirms that the gap at $\mathrm{lr}=0.01$ was an optimization issue
(insufficient learning rate for a 10-task setting with 20 classes per task),
not a fundamental limitation of the method.
Additional runs with $\mathrm{lr}\in\{0.03, 0.05\}$ are ongoing.

\textbf{Strong performance on 3 of 4 datasets.}
The partial CIFAR-100 run reaches FAA $= 88.37$ at task~7, already above the
EWC-LoRA reference of $87.91$.
The ImageNet-R partial run (task~9) gives FAA $= 77.13$ ($+4.3$pp vs.\ EWC-LoRA).
The DomainNet partial run (task~2) gives FAA $= 73.69$, matching EWC-LoRA ($73.46$).
Together with the ImageNet-A LR-tuned result, FS-LoRA is competitive or superior
on all four standard benchmarks when the LR is appropriately set.

\subsection{Hyperparameter Sensitivity}
\label{sec:hparam_main}

The hyperparameter sensitivity analysis ($\lambdaewc$ sweep on CUB200) is reported in
Appendix~\ref{sec:hparam}.
The key finding is that FS-LoRA FAA varies by less than $0.35$pp across
$\lambdaewc \in \{500, 2000, 5000\}$,
and $\cosff$ remains in $[-0.030, -0.028]$ across all values,
confirming that the orthogonality is insensitive to the Fisher penalty strength.
